\section{总结与延伸}

这堂课最重要的贡献不是宣称 AI 已经替代科学家，而是给出一套\textbf{把基础模型组织成科学协作系统}的框架。模型提供广泛知识和推理，Agent 把搜索操作分工，tournament 提供相对优先级，tools 把文本连接到专门计算，memory 保留跨轮状态，而人类负责目标、约束、实验、伦理和最终解释。

\subsection{把整堂课压缩成一个闭环}

前面各章分别讨论了搜索、排序、实验与安全；这里把它们重新连成一个有输入、有反馈、有停止条件的循环。读这条链时不要把箭头理解为一次线性流水线：实验失败会返回 Goal 与 Generate，专家也可以在任意阶段修改约束、暂停运行或要求补充证据。AI co-scientist 的工作流可以压缩为
$$
\text{Goal}
\rightarrow \text{Generate}
\rightarrow \text{Debate}
\rightarrow \text{Rank}
\rightarrow \text{Evolve}
\rightarrow \text{Test}
\rightarrow \text{Update}.
$$
其中：
\begin{itemize}
  \item Goal 由科学家给出，并可在中途重构；
  \item Generate 与 Evolve 扩大候选空间；
  \item Debate 与 Rank 暴露弱点并分配内部注意力；
  \item Test 把候选接入工具、专家和实验；
  \item Update 同时更新候选、证据图、失败记录和下一轮问题。
\end{itemize}

\begin{importantbox}{本讲的十条核心结论}
\begin{enumerate}
  \item 科学发现速度应按被验证知识衡量，而不是生成文本数量。
  \item 开放式 hypothesis generation 需要系统流程，不能只靠一次 prompt。
  \item AlphaFold 证明专门任务可以被巨大压缩，但不等于 general scientific intelligence。
  \item Multi-agent 的价值来自异质分工、比较与反馈，不来自复制相同回答。
  \item Elo 与 model-as-judge 只提供内部 prioritization，不能替代现实证据。
  \item AMR 是隐藏结果复现，不应改写成前瞻发现。
  \item AML、organoid、结构预测和课堂未发表结果属于不同证据层。
  \item Prospective validation 的关键是实验前冻结预测和指标。
  \item 长报告必须保留 provenance、反例、负结果和资源元数据。
  \item Scientist-in-the-loop 与 layered safety 都要贯穿运行过程，而不是最后补签。
\end{enumerate}
\end{importantbox}

\subsection{从 hypothesis throughput 到 validation throughput}

未来系统的限制很可能不是生成不出假设，而是没有足够实验、专家和审稿能力。工程优化应从“每小时多少候选”转向“每单位验证成本排除多少机制、获得多少信息”。这也意味着负结果、校准和停止建议与命中结果同样有价值。

\begin{knowledgebox}{部署前决策清单}
问题是否可证伪？输入资料是否冻结并可追溯？Agent 分工是否真的异质？judge rubric 是否公开？是否记录全部候选和负结果？实验前是否冻结预测？报告是否区分事实、推断和建议？是否公开 compute 区间？工具权限是否最小化？是否有人类能在运行中暂停、重构和拒绝目标？
\end{knowledgebox}

\subsection{仍然开放的研究问题}

第一，怎样评价真正 novelty，而不奖励模型熟悉的写作风格？第二，怎样让多个 Agent 共享事实又保持假设多样性？第三，怎样学习实验价值而不把昂贵现实世界反馈简化成 proxy？第四，怎样为科学报告提供 edge-level provenance 和可执行 falsification test？第五，如何让安全系统理解长时间、跨 Agent 组合后才出现的危险意图？

\begin{warningbox}{不能从本讲推出的结论}
不能推出 AI co-scientist 已能独立完成科学；不能推出所有展示案例均已发表或可临床使用；不能推出更多 test-time compute 必然产生突破；不能推出一个 10\% 监控阈值足以保障安全；也不能根据结尾过渡卡补写未录制的 AMIE 内容。
\end{warningbox}

\subsection{拓展阅读}

建议先读课堂日前的 `Towards an AI co-scientist` arXiv v1，重点检查 architecture、test-time scaling、AML、AMR 与 OCT4 appendix；再读 AMR bioRxiv v1，核查隐藏结果复现的时间线和候选比较；随后阅读 liver-fibrosis 与 plant-assembly 两篇 bioRxiv v1，比较 organoid 与结构生物学的验证层级。阅读后续 Nature 版本时，应明确它发表于课堂日期之后，不能倒推为课堂当日证据。

\begin{knowledgebox}{Lecture snapshot}
\raggedright
本讲发生于 2026-05-14。讲义使用当时已公开的四份版本：AI co-scientist arXiv v1（2502.18864）、AMR bioRxiv v1（2025.02.19.639094）、liver-fibrosis bioRxiv v1（2025.04.29.651320）与 plant-assembly bioRxiv v1（2026.05.03.722499）。2026-06-29 的 arXiv v2、后续 Nature 版本及其他课后材料不作为课堂事实。
\end{knowledgebox}

\enlargethispage{3\baselineskip}
最终，AI co-scientist 的价值不在于输出一篇“像论文”的答案，而在于把开放问题转成一组可比较、可反驳、可路由给合适专家和实验的候选。只要验证、责任和安全仍留在闭环中，Agent 才是在加速科学，而不是加速未经验证的确信。

\end{document}
